\section{Hipótesis}
\label{chap:hipotesis}

\subsection{Hipótesis de Trabajo}
\label{sec:hipotesis_trabajo}

\textbf{No aplica} según el alcance \textbf{descriptivo} adoptado en esta investigación \cite{sampieri2023metodologia}. El estudio se centra en \textbf{describir y documentar} el artefacto construido (módulo de simulación, procedimiento de cobertura y capa auxiliar MCP), sus métricas registradas y su integración en un proyecto Ionic de referencia, sin plantear una comprobación formal de relación causa--efecto entre variables independientes y dependientes.

La evidencia se presentará como \textbf{caracterización y registro medible} (cobertura de líneas y ramas, tiempos de preparación de pruebas, pasaje de suite, lecturas opcionales de calidad estática) tomada en un momento definido del desarrollo. Si en etapas posteriores el asesor y el jurado consideran pertinente elevar el alcance a correlacional o explicativo, se reformulará la hipótesis en la memoria final con el diseño experimental correspondiente.
